\chapter{Tablice}\markboth{Tablice}{Tablice}

Tablice w tym dodatku są generowane z dwóch plików danych (tempo i długości książek)
przez \code{tools/gen.py}. Zmiana założenia zmienia więc tabelę,
a nie tylko tekst obok niej.

\section{Wzór}
Niech $m$ oznacza liczbę minut czytania dziennie, $v$ tempo w słowach na minutę, $w$ długość
książki w słowach, a $N$ liczbę książek w roku. Wtedy
\[
N = \frac{365\,m\,v}{w}
\qquad\text{oraz}\qquad
m = \frac{N\,w}{365\,v}.
\]
Pierwszy wzór czyta się tak: dzienny czas razy liczba dni daje minuty w roku; minuty razy tempo
daje słowa w roku; słowa w roku podzielone przez długość książki dają liczbę książek. Sprawdzenie
dla $m = 10$, $v = 150$ i $w = 80\,000$: $3650 \cdot 150 = 547\,500$ słów, a to dzieli się na
$80\,000$ z resztą, czyli około $6{,}8$ książki, co zgadza się z tabelą poniżej.

Założenia tej tabeli są następujące:
\begin{itemize}
\item typowa książka ma 80\,000 słów, długa 120\,000, krótka 40\,000;
\item tempo jest stałe, bez przerw, bez cofania się i bez spadku uwagi;
\item rok ma 365 dni, każdego dnia czytasz, bez wolnych;
\item nie liczymy czasu na wybór książki, notatki, powroty do wcześniejszych stron ani zmianę książki;
\item tempo 250 słów na minutę jest założeniem modelu, a nie wynikiem pomiaru.
\end{itemize}
Trzy z tych założeń (stałe tempo, czytanie każdego dnia i pominięcie czasu na wybór i notatki)
w prawdziwym życiu działają w jedną stronę: zawyżają liczbę książek. Tabela jest więc raczej górną
granicą tego, co arytmetyka pozwala zrobić, a nie prognozą.

\section{Książki w roku}
Poniższa tabela podaje liczbę typowych książek (80\,000 słów) w roku. Wiersze to minuty czytania
dziennie, kolumny to tempo w słowach na minutę.

\generated{books-per-year}

\section{Długość książki w godzinach}
Poniżej czas czytania trzech długości książek przy 250 słowach na minutę. Dla innego tempa
dziel słowa przez tempo i przelicz na godziny.

\generated{book-lengths}

\section{Jak policzyć własną tablicę}
Do własnej tablicy potrzebne są dwie liczby: tempo i długość książki. Oba trzeba zmierzyć, a nie
przyjąć z tabeli.

\subsection*{Krok 1: tempo}
Weź trzy strony tekstu, który już znasz i który jest mniej więcej tak trudny, jak lektury, które
planujesz. Policz słowa na stronie: na jednej stronie wystarczy policzyć słowa w kilku wierszach
i pomnożyć przez liczbę wierszy. Zmierz czas czytania zegarem, bez czytania na głos. Podziel
słowa przez minuty. Przykład: trzy strony po 300 słów to 900 słów; przeczytane w 4,5 minuty
dają tempo $900 / 4{,}5 = 200$ słów na minutę. Zmierz tak dwa razy w różnych dniach i weź wynik
bliższy średniej, a nie najlepszy.

\subsection*{Krok 2: długość książki}
Strony wydawcy trzeba przeliczyć na słowa, a stosunek słów do stron zależy od czcionki, marginesów
i tego, czy w książce są rysunki. Zmierz go na jednej stronie tej samej książki, zanim przyjmiesz
go dla całej.

\subsection*{Krok 3: cel w książkach}
Zapisz liczbę książek, którą chcesz przeczytać w roku, i przeciętną długość tych, które naprawdę
planujesz. Nie zapisuj długości z pamięci, bo zwykle pamiętamy krótkie tytuły.

\subsection*{Krok 4: podstawienie}
Przykład: tempo 200 słów na minutę, cel 20 książek po 60\,000 słów. Wtedy
\[
m = \frac{20 \cdot 60\,000}{365 \cdot 200} = \frac{1\,200\,000}{73\,000} \approx 16{,}4
\]
minuty dziennie. Przy tempie 250 słów na minutę ta sama liczba książek wymagałaby około
13 minut dziennie. Różnica w tempie jest tu tak samo ważna jak wybór długości książek.

\begin{falle}
Najczęstszy błąd to przyjęcie tempa z internetowej tabeli zamiast z własnego pomiaru. Tabela
mierzy kogoś innego, a licznik pokazuje dokładnie to, co wstawisz. Jeśli wpiszesz cudze tempo,
dostaniesz cudzy rok.
\end{falle}
